\documentclass[12pt,a4paper]{article}

% --- Pacotes Básicos ---
\usepackage[utf8]{inputenc}
\usepackage[english]{babel}
\usepackage[T1]{fontenc}
\usepackage{graphicx}
\usepackage{amsmath, amssymb}
\usepackage{booktabs}
\usepackage{geometry}
\usepackage{hyperref}
\usepackage{setspace}

\geometry{left=2cm, right=2cm, top=3cm, bottom=2cm}
\onehalfspacing

\title{K-Nearest Neighbour Modification and Performance Benchmarking}
\author{
    Filipe Zheng (202406753) \\
    Maiara Sol Ribeiro Guedes (202407694) \\
    Sílvia Cármen Baciu Pinto (202405988)
}
\date{\today}

\begin{document}
\maketitle
\newpage

\section*{Introduction}
\addcontentsline{toc}{section}{Introduction}
\begin{itemize}
    \item \textbf{Goals:} []
    \item \textbf{Approach:} []
    \item \textbf{Results:} []
\end{itemize}
\newpage

% --- Índice ---
\tableofcontents
\newpage

% --- Introdução e Algoritmo Selecionado ---
\section{Selecting the Algorithm}
\subsection{K-Nearest Neighbour Implementation}
% Descreva aqui o funcionamento lógico do algoritmo base (ex: SVM, Random Forest).

\subsection{Initial Performance Benchmarking}
% Discuta como o algoritmo se comporta com ruído, outliers ou classes desequilibradas.

% --- Proposta ---
\section{Algorithm Modification}
\subsection{Reasons}
% Por que é necessária uma alteração? Qual o problema identificado no algoritmo base?

\subsection{Description}
% Detalhe a alteração técnica, fórmulas ou mudanças na lógica do classificador.

% --- Estudo Empírico ---
\section{Empirical Study}
\subsection{Testing methodology}
% Metodologia de teste (ex: Cross-validation).

\subsection{Datasets}
% Tabela com as características dos dados (instâncias, atributos, rácio de desequilíbrio).
\begin{table}[h!]
\centering
\begin{tabular}{@{}llll@{}}
\toprule
Dataset & Instances & Attributes & Variable of interest \\ \midrule
Data 1  & 1000       & 20        & Imbalance (1:10)            \\
Data 2  & 500        & 10        & Noise (5\%)                 \\ \bottomrule
\end{tabular}
\caption{Datasets}
\end{table}

\subsection{Hiperparameters}
% Lista de parâmetros afinados durante as experiências.

\subsection{Performance Benchmarking}
% Métricas utilizadas (Accuracy, F1-Score, AUC).

% --- Resultados ---
\section{Final Analysis}
\subsection{Results}
% Gráficos e tabelas comparativas entre o algoritmo original e o proposto.

\subsection{Discussion}
% Interpretação dos dados à luz das hipóteses levantadas.

% --- Conclusão ---
\section{Conclusion}
% Resumo das contribuições e limitações, com sugestões para investigação posterior.

% --- Referências ---
\begin{thebibliography}{9}
\bibitem{ref1} Gama, J., Carvalho A., Faceli K., Lorena, A.C. and Oliveira, M., \textit{Extração de Conhecimento de Dados: Data Mining}, Edições Sílabo, 2017
\bibitem{ref2} J. Moreira, A. Carvalho, and T. Horvath, \textit{A General Introduction to Data Analytics}, John Wiley \& Sons, 2018
\end{thebibliography}

\end{document}
